% Faithful translation baseline; Chinese source PDF pp. 41--42.
\section{Coherent States of the Harmonic Oscillator}
In energy eigenstates of the harmonic oscillator, the expectation values of $\op x$ and $\op p$ are zero, and no classical oscillation is observed. However, consider a linear combination of energy eigenstates, for example
\begin{equation}
\ket a=c_0\ket0+c_1\ket1.
\end{equation}
In the Heisenberg picture, the expectation value is
\begin{equation}
\begin{split}
\mel a{\op x(t)}a
&=(c_0^*\bra0+c_1^*\bra1)\op x(t)(c_0\ket0+c_1\ket1)\\
&=c_0^*c_1\mel0{\op x(t)}1+c_1^*c_0\mel1{\op x(t)}0\\
&=\sqrt{\frac\hbar{2m\omega}}
\left(c_0^*c_1e^{-\ii\omega t}+c_1^*c_0e^{\ii\omega t}\right).
\end{split}
\end{equation}
Oscillations appear. Can one find a state whose oscillation is closest to classical oscillation, that is, the oscillation of $\op x(t)$? This leads to the coherent state $\ket\alpha$, defined as an eigenstate of the lowering operator:
\begin{equation}
\boxed{\op a\ket\alpha=\alpha\ket\alpha.}
\end{equation}
Note that $\op a$ is not Hermitian; its eigenvalue $\alpha$ is generally complex.

\subsection*{1) Coherent states satisfy the minimum-uncertainty relation}
The uncertainty relation is
\begin{equation}
(\Delta x)^2(\Delta p)^2\ge\frac{\hbar^2}{4}.
\end{equation}
In the coherent state $\ket\alpha$,
\begin{equation}
\begin{split}
\mel\alpha{\op x}\alpha
&=\sqrt{\frac\hbar{2m\omega}}\mel\alpha{\op a^{\dagger}+\op a}\alpha
=\sqrt{\frac\hbar{2m\omega}}(\alpha^*+\alpha),\\
\mel\alpha{\op x^2}\alpha
&=\frac\hbar{2m\omega}\bigl[1+(\alpha+\alpha^*)^2\bigr],\\
\mel\alpha{\op p}\alpha
&=-\ii\sqrt{\frac{\hbar m\omega}{2}}(\alpha-\alpha^*),\\
\mel\alpha{\op p^2}\alpha
&=\frac{\hbar m\omega}{2}\bigl[1-(\alpha-\alpha^*)^2\bigr].
\end{split}
\end{equation}
It follows that
\begin{equation}
(\Delta x)^2(\Delta p)^2
=\bigl(\langle\op x^2\rangle-\langle\op x\rangle^2\bigr)
\bigl(\langle\op p^2\rangle-\langle\op p\rangle^2\bigr)
=\frac{\hbar^2}{4}.
\end{equation}
Thus coherent states satisfy the minimum-uncertainty relation and are the quantum states closest to classical states.

\subsection*{2) Expansion of a coherent state}
Expand $\ket\alpha$ in energy eigenstates:
\begin{equation}
\ket\alpha=\sum_n c_n\ket n.
\end{equation}
Since
\begin{equation}
c_n=\braket n\alpha
=\frac1{\sqrt{n!}}\bra0\op a^n\ket\alpha
=\frac{\alpha^n}{\sqrt{n!}}\braket0\alpha
=\frac{\alpha^n}{\sqrt{n!}}c_0,
\end{equation}
we have
\begin{equation}
\ket\alpha=c_0\sum_n\frac{\alpha^n}{\sqrt{n!}}\ket n.
\end{equation}
Normalization gives
\begin{equation}
1=\braket\alpha\alpha=\sum_n|c_n|^2
=|c_0|^2\sum_n\frac{|\alpha|^{2n}}{n!}
=|c_0|^2e^{|\alpha|^2}.
\end{equation}
Choose $c_0=e^{-|\alpha|^2/2}$; then
\begin{equation}
\boxed{\ket\alpha=\sum_{n=0}^{\infty}
e^{-|\alpha|^2/2}\frac{\alpha^n}{\sqrt{n!}}\ket n.}
\end{equation}

\subsection*{3) The ground state is a coherent state}
For the ground state $\ket0$, $\op a\ket0=0$. Thus $\ket0$ is an eigenstate of the annihilation operator, with eigenvalue zero.

\subsection*{4) Time evolution of the eigenvalue}
In the Heisenberg picture,
\begin{equation}
\op a(t)=\op a(0)e^{-\ii\omega t},
\end{equation}
\begin{equation}
\op a(t)\ket\alpha=\op a(0)e^{-\ii\omega t}\ket\alpha
=\alpha e^{-\ii\omega t}\ket\alpha.
\end{equation}
Hence the eigenvalue evolves as
\begin{equation}
\boxed{\alpha(t)=\alpha e^{-\ii\omega t}.}
\end{equation}

\interactiveguide{IC-02-03}{Coherent-state dynamics}{After obtaining $alpha(t)$, prepare an oscillator coherent state and vary its amplitude and phase. Compare the Fock distribution and the moving position density; revisit the experiment after studying squeezed states to see spreading.}
